\chapter{Geradores de analisadores sintáticos}

\section{Introdução}

No capítulo anterior estudamos os algoritmos LR(1) e LALR(1) de forma
teórica e implementamos suas tabelas em Haskell. Na prática, contudo,
construir essas tabelas à mão para uma gramática real seria
impraticável. Assim como o Alex automatiza a construção de
analisadores léxicos a partir de expressões regulares, o
\textbf{Happy} automatiza a construção de analisadores sintáticos a
partir de especificações gramaticais no estilo BNF. O Happy gera
analisadores LALR(1) e é a ferramenta padrão para este fim no
ecossistema Haskell: o próprio compilador GHC usa um parser gerado
pelo Happy.

Este capítulo apresenta o Happy através da especificação do parser de
expressões aritméticas, disponível em
\begin{flushleft}
  \verb|Exp/Frontend/Parser/Happy/ExpParser.y|
\end{flushleft}
Mostramos como cada parte do arquiv \texttt{.y} se
relaciona com os conceitos vistos nos capítulos anteriores e como
usar a opção \verb|-i| para inspecionar a tabela
LALR gerada. Por fim, ilustramos os dois tipos de conflito que podem
surgir em gramáticas ambíguas: conflitos shift-reduce e conflitos reduce-reduce.

\section{Estrutura de uma especificação Happy}

Um arquivo de especificação Happy (extensão .y) é dividido em quatro
partes, separadas
pelo marcador \%\%:

\begin{verbatim}
{ cabeçalho Haskell }

diretivas e declaração de tokens

%%

regras gramaticais

{ rodapé Haskell }
\end{verbatim}

\subsection{Cabeçalho}

A primeira seção, delimitada por chaves, contém código Haskell
puro que é copiado literalmente para
o início do arquivo gerado. É aqui que se declara o nome do módulo e
os imports necessários:

\begin{minted}{haskell}
{
module Exp.Frontend.Parser.Happy.ExpParser (expParser, parserTest) where

import Exp.Frontend.Lexer.Token
import Exp.Frontend.Lexer.Alex.ExpLexer hiding (lexer)
import Exp.Frontend.Syntax.ExpSyntax
}
\end{minted}

\subsection{Diretivas}

As diretivas configuram o parser gerado. As mais comuns são:

\begin{table}[H]
  \begin{tabular}{p{0.5\linewidth} p{0.5\linewidth}}
    \toprule
    Diretiva & Significado \\
    \midrule
    \verb|%name p N| & Gera a função \verb|p| que inicia o parsing
    pelo não-terminal \verb|N| \\
    \verb|%tokentype { T }| & Define o tipo Haskell dos tokens consumidos \\
    \verb|%monad { M }{ (>>=) }{ return }| & Executa o parser dentro
    da mônada \verb|M| \\
    \verb|%lexer { f }{ eof }| & Integra um lexer incremental;
    \verb|f| obtém o próximo token e \verb|eof|
    é o token de fim de entrada \\
    \verb|%error { h }| & Chama a função \verb|h| quando ocorre um
    erro de sintaxe \\
    \verb|%expect n| & Declara que a gramática possui exatamente
    \verb|n| conflitos shift-reduce esperados; o Happy falha se o
    número for diferente \\
    \bottomrule
  \end{tabular}
  \centering
\end{table}

No parser de expressões, essas diretivas ficam assim:

\begin{verbatim}
%name parser Exp
%monad {Alex}{(>>=)}{return}
%tokentype { Token }
%error     { parseError }
%lexer {lexer}{Token _ TEOF}
\end{verbatim}

A diretiva \verb|%monad| integra o Happy com a
mônada \verb|Alex| do lexer gerado pelo Alex. Com
\verb|%lexer|, o Happy chama a função
\haskell{lexer} cada vez que precisa de um token, em
vez de receber a lista completa de antemão. Esse modo incremental é
necessário quando o lexer é monádico.

\subsection{Declaração de tokens}

A seção \verb|%token| mapeia nomes simbólicos
usados nas regras gramaticais para padrões de correspondência (em
Haskell) contra o tipo de token. Cada linha tem a forma:

\begin{verbatim}
nome_simbólico  { padrão Haskell }
\end{verbatim}

No parser de expressões:

\begin{verbatim}
%token
      num       {Token _ (TNumber $$)}
      '('       {Token _ TLParen}
      ')'       {Token _ TRParen}
      '+'       {Token _ TPlus}
      '*'       {Token _ TTimes}
\end{verbatim}

O uso de \verb|$$| no padrão
\haskell{Token \_ (TNumber $$)} é especial: ele
captura o valor inteiro contido no token e o torna disponível nas
ações semânticas como \verb|$1|,
\verb|$2|, etc. Para tokens sem valor associado,
como \verb|TLParen|, o padrão apenas verifica que
o construtor bate.

\subsection{Declarações de precedência}

As declarações \verb|%left|,
\verb|%right| e
\verb|%nonassoc| servem para resolver conflitos
shift-reduce que surgem em gramáticas de expressões escritas de forma
ambígua. Cada linha associa uma precedência (determinada pela ordem
das declarações, de menor para maior) e uma associatividade a um ou mais tokens:

\begin{verbatim}
%left '+'
%left '*'
\end{verbatim}

Neste exemplo, \verb|*| tem precedência mais alta
que \verb|+| (pois aparece depois), e ambos são
associativos à esquerda. Veremos na seção sobre conflitos como essas
declarações eliminam a ambiguidade.

\subsection{Regras gramaticais}

Após o marcador \verb|%%|, definem-se as
produções da gramática no estilo BNF aumentado com ações semânticas em Haskell:

\begin{verbatim}
Exp : num         { EInt $1 }
    | Exp '+' Exp { $1 :+: $3 }
    | Exp '*' Exp { $1 :*: $3 }
    | '(' Exp ')' { $2 }
\end{verbatim}

Cada alternativa é uma produção seguida de uma ação entre chaves. Os
símbolos \verb|$1|, \verb|$2|,
\verb|$3| referem-se, respectivamente, ao
primeiro, segundo e terceiro símbolo do lado direito da produção. A
ação é executada quando o parser reduz por aquela produção,
construindo o nó correspondente da AST.

O tipo de retorno de cada não-terminal é inferido pelo compilador
Haskell a partir das ações: como \haskell{EInt $1} e
\haskell{$1 :+: $3} têm tipo
\haskell{Exp}, o não-terminal
\haskell{Exp} produz valores do tipo
\haskell{Exp}.

\subsection{Rodapé}

A última seção, também delimitada por chaves, contém funções
auxiliares usadas nas ações ou nas diretivas:

\begin{minted}{haskell}
{
parserTest :: String -> IO ()
parserTest s = do
  r <- expParser s
  print r

parseError (Token (line, col) lexeme)
  = alexError $ "Parse error while processing lexeme: " ++ show lexeme
                ++ "\n at line " ++ show line ++ ", column " ++ show col

lexer :: (Token -> Alex a) -> Alex a
lexer = (=<< alexMonadScan)

expParser :: String -> IO (Either String Exp)
expParser content = do
  pure $ runAlex content parser
}
\end{minted}

A função \haskell{parseError} recebe o token que
causou o erro e usa \haskell{alexError} para propagar
a mensagem de erro dentro da mônada Alex. A função
\haskell{lexer} adapta
\haskell{alexMonadScan} (que produz
\haskell{Alex Token}) para a interface esperada pelo
Happy (que recebe uma continuação
\haskell{Token -> Alex a}).

\section{Gerando o parser e inspecionando a tabela}

Para gerar o parser a partir da especificação, basta invocar o Happy:

\begin{minted}{bash}
happy ExpParser.y
happy -o ExpParser.hs ExpParser.y
\end{minted}

Quando a especificação faz parte de um projeto Cabal, o próprio Cabal
detecta arquivos .y em
\verb|build-tool-depends| (com
happy listado) e os processa
automaticamente antes da compilação Haskell.

\subsection{A opção -i}

A opção \verb|-i| instrui o Happy a gerar um
\textbf{arquivo de informação} (por padrão
\verb|ExpParser.info|) que descreve em detalhe a
tabela LALR construída:

\begin{minted}{bash}
happy -i ExpParser.y
\end{minted}

O arquivo gerado contém:

\begin{enumerate}
  \item
    \textbf{A gramática aumentada}: as produções numeradas,
    incluindo a produção inicial $S' \to Exp$.
  \item
    \textbf{Conjuntos FIRST} de cada não-terminal.
  \item
    \textbf{Os estados do autômato}: cada estado lista seus itens
    LR(1), com os lookaheads calculados pelo algoritmo LALR.
  \item
    \textbf{A tabela de ações}: para cada par (estado, terminal),
    a ação: shift, reduce ou accept.
  \item
    \textbf{A tabela goto}: para cada par (estado, não-terminal),
    o estado destino.
  \item
    \textbf{Conflitos}: se houver, eles são listados com o estado
    e o token envolvidos, indicando qual ação o Happy escolheu (por
    padrão, shift vence em conflitos shift-reduce).
\end{enumerate}

Este arquivo é extremamente útil para depurar gramáticas: ao ver em
qual estado um conflito ocorre e quais itens ele contém, é possível
identificar a ambiguidade na gramática e corrigi-la com declarações
de precedência ou reestruturando as produções.

\section{Conflitos shift-reduce}

Um conflito shift-reduce ocorre quando, em um determinado estado, o
parser pode tanto deslocar (shift) o próximo token quanto reduzir por
alguma produção. O caso mais comum é o de gramáticas de expressões
escritas de forma ambígua.

Considere a gramática de expressões \textbf{sem} declarações de precedência:

\begin{verbatim}
Exp : num
    | Exp '+' Exp
    | Exp '*' Exp
    | '(' Exp ')'
\end{verbatim}

Ao analisar a entrada \verb|1 + 2 * 3|, após ter
empilhado \verb|Exp '+' Exp| (com o segundo
\verb|Exp| valendo 2), o
parser vê * na entrada. Ele pode:

\begin{itemize}
  \item
    \textbf{Reduzir} \verb|Exp '+' Exp| para
    \verb|Exp| (interpretando
    \verb|(1 + 2) * 3|), ou
  \item
    \textbf{Deslocar} \verb|*| (interpretando
    \verb|1 + (2 * 3)|).
\end{itemize}

Sem orientação adicional, o Happy reporta o conflito e, por padrão,
escolhe sempre o shift. A mensagem no terminal seria:

\begin{verbatim}
ExpParser.y: shift/reduce conflicts:  4
\end{verbatim}

E no arquivo .info, cada conflito é
detalhado, por exemplo:

\begin{verbatim}
State 7

        Exp -> Exp . '+' Exp          (rule 2)
        Exp -> Exp . '*' Exp          (rule 3)
        Exp -> Exp '+' Exp .          (rule 2)

        '*'    shift, and enter state 5
               (reduce using rule 2)

        '+'    shift, and enter state 4
               (reduce using rule 2)

        ')'    reduce using rule 2
        '$'    reduce using rule 2
\end{verbatim}

A resolução correta é declarar precedência e associatividade antes do
\%\%. Ao escrevermos:

\begin{verbatim}
%left '+'
%left '*'
\end{verbatim}

o Happy usa essas informações para resolver cada conflito: quando o
token a deslocar tem precedência maior que a produção a reduzir, faz
shift; quando tem menor, reduz; quando são iguais, a associatividade
decide. Com as declarações acima, \verb|*| tem
precedência maior que \verb|+|, portanto
\verb|1 + 2 * 3| é interpretado como
\verb|1 + (2 * 3)|, e
\verb|%expect 0| pode ser declarado para
confirmar que nenhum conflito resta.

Outro exemplo clássico de conflito shift-reduce é o do \textbf{else
pendente} (\emph{dangling else}), presente em linguagens com
construção \verb|if-then-else| opcional:

\begin{verbatim}
Stmt : 'if' Exp 'then' Stmt
     | 'if' Exp 'then' Stmt 'else' Stmt
     | 'other'
\end{verbatim}

Ao analisar if e1 then if e2 then s1 else s2, após empilhar o
then interior, o parser vê else e tem dúvida: este
else pertence ao
if externo ou ao interno? Shift associa o
else ao if mais
próximo (interno), que é o comportamento convencional na maioria das
linguagens. A declaração \verb|%shift!| ou
simplesmente aceitar o comportamento padrão do Happy resolve o conflito.

\section{Conflitos reduce-reduce}

Um conflito reduce-reduce ocorre quando, em um determinado estado e
para um mesmo token de lookahead, dois itens completos diferentes
indicam reduções por produções distintas. Isso geralmente revela uma
ambiguidade estrutural na gramática que não pode ser resolvida por precedência.

Considere a gramática:

\begin{verbatim}
Expr : termo
     | termo

termo : 'x'
      | 'x'
\end{verbatim}

De forma mais realista, o conflito aparece quando dois não-terminais
distintos derivam a mesma sequência de terminais e são
intercambiáveis no mesmo contexto:

\begin{verbatim}
Prog : Decl
     | Expr

Decl : 'id'
Expr : 'id'
\end{verbatim}

Ao ver o token \verb|id| e precisar reduzir, o
parser não sabe se deve criar um nó \verb|Decl| ou
um nó \verb|Expr|. O arquivo
.info mostraria algo como:

\begin{verbatim}
State 1

        Decl -> 'id' .          (rule 1)
        Expr -> 'id' .          (rule 2)

        '$'    reduce using rule 1
               (reduce using rule 2)
\end{verbatim}

O Happy resolve o conflito escolhendo sempre a \textbf{primeira
regra} na ordem em que aparece no arquivo, mas emite um aviso.
Diferentemente dos conflitos shift-reduce, conflitos reduce-reduce
raramente têm uma resolução correta por padrão: eles quase sempre
indicam que a gramática precisa ser reformulada.

A solução mais comum é tornar os dois não-terminais sintáticos
distintos por contexto, introduzindo palavras-chave discriminadoras:

\begin{verbatim}
Prog : 'let' 'id'        -- declaração
     | 'id'              -- expressão
\end{verbatim}

Ou, alternativamente, desambiguar na fase semântica (análise de tipos
ou tabela de símbolos), como fazem as gramáticas de linguagens como
C, onde typedef torna certos
identificadores indistinguíveis de tipos sem informação de contexto.

\section{Conclusão}

Este capítulo apresentou o Happy, o gerador de analisadores
sintáticos padrão em Haskell. Vimos como um arquivo
.y é estruturado em quatro partes
(cabeçalho, diretivas, regras e rodapé) e como cada diretiva
configura o parser gerado. Estudamos a especificação concreta do
parser de expressões aritméticas, que integra o Happy com o lexer
monádico gerado pelo Alex. Mostramos como a opção
-i permite inspecionar os estados LALR e os
conflitos detectados, tornando o processo de depuração de gramáticas
transparente. Por fim, ilustramos os dois tipos de conflito
(shift-reduce, geralmente resolúvel com declarações de precedência, e
reduce-reduce, que normalmente exige a reformulação da gramática),
conectando os conceitos teóricos dos capítulos anteriores à prática
do desenvolvimento de compiladores em Haskell.
